\documentclass[11pt]{article}
\usepackage[utf8]{inputenc}
\usepackage[T1]{fontenc}
\usepackage{lmodern}
\usepackage{geometry}
\usepackage{hyperref}
\usepackage{enumitem}
\usepackage{graphicx}
\usepackage{xcolor}
\geometry{margin=1in}

\title{Prompt Engineering Midterm Examination Assignment}
\author{Dr.~Mehmet Ali Akyol}
\date{Fall 2025}

\begin{document}
\maketitle

\section*{Purpose and Scope}
This midterm examination replaces an in-class test and evaluates your mastery of the core course pillars covered to date: foundations of prompt design, model mechanics, reusable patterns and frameworks, tone/format control, and iterative refinement. The assignment is deliberately structured to require authentic experimentation, documentation, and reflection. Merely querying a large language model (LLM) for a finished answer will not satisfy the requirements. You are expected to:
\begin{itemize}
    \item Design and iterate on prompts across three distinct domains.
    \item Collect human feedback, observational notes, and supporting evidence.
    \item Critically evaluate outputs using rubrics you construct.
    \item Explain your reasoning with references to course concepts and activities.
\end{itemize}

\section*{Honor Statement}
Submit a signed statement (typed name acceptable) confirming that every artifact is your original work, created expressly for this exam, and that you complied with the AI usage policy: LLMs may assist only as tools within the documented workflow. No LLM-generated text should appear unreviewed or unannotated.

\section{Deliverables Overview}
Submit a single PDF dossier (8--12 pages) containing the following sections. Supplemental files (recordings, tables, raw data) must be placed in an appendix folder and linked from the PDF.

\subsection{Section~1: Scenario Portfolio}
\begin{enumerate}[label=1.\arabic*]
    \item Select \textbf{three domains} relevant to your academic or professional interests (e.g., healthcare, legal, education, finance, engineering). Provide a one-paragraph justification for each choice referencing our early discussions on audience, purpose, and constraints.
    \item For each domain, articulate a real-world task and describe the stakeholders and success criteria. Use the prompt brief template introduced in Week~1.
    \item Include at least one artifact (photo of whiteboard, scanned notes, stakeholder interview summary) evidencing offline ideation. Annotate with timestamps and context.
\end{enumerate}

\subsection{Section~2: Model Mechanics Insight Log}
\begin{enumerate}[label=2.\arabic*]
    \item Create an \textbf{experiment log} analyzing how token limits, temperature, and few-shot examples affect the same task across two models of different scales. Provide:
    \begin{itemize}
        \item Screenshots or copy/pasted raw outputs with token counts visible.
        \item A short calculation estimating token usage and cost.
    \item Commentary explaining observed differences using course terminology (tokens, embeddings, self-attention, emergent behavior).
    \end{itemize}
    \item Include a hand-drawn or self-created diagram (no autogenerated graphics) illustrating your mental model of the transformer attention flow for the chosen task.
\end{enumerate}

\subsection{Section~3: Prompt Pattern Stack}
\begin{enumerate}[label=3.\arabic*]
    \item For each domain task, design a \textbf{pattern stack} incorporating at least three atomic patterns (role, format, reasoning, context, constraints) and one named framework introduced in class (e.g., CO-STAR, IDEA Loop). Present:
    \begin{itemize}
        \item Initial prompt draft (v1).
        \item Pattern stack rationale table mapping instructions to expected benefits.
    \item A critique of v1 using a rubric you author.
    \item Revised prompt (v2) with annotations showing how feedback was addressed.
    \end{itemize}
\end{enumerate}

\subsection{Section~4: Tone, Style, and Format Stress Test}
\begin{enumerate}[label=4.\arabic*]
    \item Construct a \textbf{tone ladder} (from formal to conversational) for one domain task. For each rung, include:
    \begin{itemize}
        \item The prompt snippet enforcing tone.
        \item Model output (truncated to 150 words) marked up with manual annotations verifying compliance.
        \item Commentary on any drift and corrective revisions.
    \end{itemize}
    \item Develop a format schema (table, JSON, or custom template) and demonstrate compliance across two models. Highlight where models failed and document your manual fixes.
\end{enumerate}

\subsection{Section~5: Iterative Refinement Case Study}
\begin{enumerate}[label=5.\arabic*]
    \item Choose the \textbf{most challenging domain task} and document three full iteration cycles:
    \begin{itemize}
        \item Cycle log showing prompt version, model configuration, and output snippet.
        \item Critique prompts (self-review or peer feedback) with scoring rubric.
    \item Revision notes referencing specific heuristics discussed in class (e.g., root-cause analysis, critique archetypes, judge models).
    \end{itemize}
    \item Produce a comparative analysis chart summarizing improvements (accuracy, completeness, tone, structure) per iteration. This chart must be built manually in a spreadsheet or plotting tool and inserted as a figure.
    \item Include a 300-word reflection on what iteration tactics were most effective and how you would automate parts of the loop responsibly.
\end{enumerate}

\section{Authenticity Requirements}
\begin{itemize}
    \item \textbf{Human Evidence:} At least two artifacts must be clearly non-digital (scanned notes, photographed sketches, sticky notes). Include date and context captions.
    \item \textbf{Peer Interaction:} Conduct a 10-minute feedback session with a classmate or colleague on one domain prompt. Attach a signed feedback form summarizing their observations and how you integrated them.
    \item \textbf{Asset Integrity:} Record all LLM interactions in a separate appendix with timestamps. Highlight where you edited or rejected model outputs.
    \item \textbf{Reflection:} Tie every major decision to specific lectures, activities, or readings from our sessions so far. Cite slide titles or page numbers where applicable.
\end{itemize}

\section{Submission Format}
\begin{itemize}
    \item Primary dossier: PDF named \texttt{MIDTERM\_LastnameFirstname.pdf}.
    \item Appendix folder (zipped): includes raw logs, media files, spreadsheets. Name \texttt{MIDTERM\_LastnameFirstname\_APPENDIX.zip}.
    \item Upload both files to the LMS by the deadline (11~November~2025, 23:59 TRT).
    \item Late submissions incur 10\% penalty per 24-hour period.
\end{itemize}

\section{Evaluation Rubric (100 points)}
\begin{itemize}[leftmargin=*]
    \item \textbf{Scenario Rigor (15)} --- Clear articulation of audience, constraints, and success criteria; tangible offline ideation evidence.
    \item \textbf{Model Mechanics Insight (15)} --- Depth of experimentation, correctness of token/cost reasoning, clarity of diagrams.
    \item \textbf{Pattern and Framework Application (20)} --- Quality of pattern stacks, critique quality, and integration of course concepts.
    \item \textbf{Tone and Format Mastery (15)} --- Effectiveness of tone ladder, schema compliance analysis, handling of drift.
    \item \textbf{Iteration Excellence (20)} --- Documentation of cycles, use of heuristics, measurable improvements, reflection depth.
    \item \textbf{Authenticity and Governance (10)} --- Completeness of artifacts, peer feedback integration, honor compliance.
    \item \textbf{Communication Quality (5)} --- Organization, clarity, visual presentation, adherence to submission guidelines.
\end{itemize}

\section*{Optional Bonus (up to 5 points)}
Design a one-page internal memo proposing how your organization (or hypothetical team) could operationalize the most successful prompt workflow from this exam. Include roles, tooling, risks, and success metrics. The memo must be written without any LLM drafting assistance (state how you ensured this).

\vfill
\noindent\textbf{Reminder:} This midterm is an individual assessment. Unauthorized collaboration or submission of AI-generated content without attribution violates the course policy and university regulations.

\end{document}
